% Source: https://artofproblemsolving.com/community/c3h2419639_how_to_tell_if_a_number_is_invertible_mod_n
% Stable identity: aops-topic-2419639-q01
\textbf{(a)} Suppose $ax\equiv1\pmod n$, so $ax-1=ny$ for some integer $y$. Any common divisor $d$ of $a$ and $n$ divides $ax-ny=1$, so $\gcd(a,n)=1$.

Conversely, suppose $\gcd(a,n)=1$. Look at the $n$ numbers $a\cdot0,\ a\cdot1,\ \dots,\ a\cdot(n-1)$. If two of them were congruent, $ai\equiv aj\pmod n$ with $0\le j<i\le n-1$, then $n\mid a(i-j)$. Since $a$ and $n$ have no common prime factor, every prime power dividing $n$ must divide $i-j$, so $n\mid i-j$, which is impossible because $0<i-j<n$. So the $n$ numbers have $n$ different remainders modulo $n$; that is, every remainder $0,1,\dots,n-1$ occurs exactly once. In particular $ax\equiv1\pmod n$ for some $x$. $\blacksquare$

\textbf{(b)} $\gcd(9,26)=1$, so $9$ is invertible modulo $26$; since $9\cdot3=27\equiv1\pmod{26}$, an inverse is $3$. But $\gcd(13,26)=13$, so $13$ is not invertible: indeed $13x$ is always $\equiv0$ or $\equiv13\pmod{26}$, never $\equiv1$.

So $\boxed{9\text{ is invertible (inverse }3\text{), }13\text{ is not}}$.
